\section{Related work}
\label{sec:related}

\paragraph{Multi-turn degradation.}
Several benchmarks evaluate conversational ability over multiple exchanges,
including MT-Bench \citep{zheng2023judging}, MT-Eval \citep{kwan2024mteval} and
MINT \citep{wang2024mint}.

\citet{laban2025lost} showed that LLMs lose substantial performance when a fully
specified instruction is sharded across turns, mainly through increased
unreliability. \citet{huang2026benefit} compared full conversational context with
contexts from which assistant turns are omitted or condensed, using in-the-wild
conversations and the Lost in Conversation math tasks. They found that response
quality is largely preserved under aggressive filtering, and that prior assistant
replies can propagate errors (context pollution).

Our Study 1 replicates this assistant-omission contrast with a frozen,
task-paired analysis. Study 2 extends the question to 100-turn dialogues with
exact answers and controlled difficulty.

Related single-turn failure modes include:
\begin{itemize}
  \item sensitivity to long inputs \citep{liu2024lost};
  \item distraction by irrelevant context \citep{shi2023distracted};
  \item unreliable self-correction \citep{huang2024cannot}.
\end{itemize}

\paragraph{Mitigation by context management.}
Two lines of work restructure the context to recover performance:
\begin{itemize}
  \item ERGO \citep{khalid2025ergo} detects entropy spikes and consolidates the
    prompt;
  \item \citet{huang2026benefit} filter assistant turns.
\end{itemize}
Our Study 2 control arm is a verified state summary. It belongs to this family,
but it uses solver-verified state, so it serves as an upper bound on what state
compression can achieve.

\paragraph{Reasoning-quality dimensions.}
Beyond final accuracy, reasoning quality has been characterized in several ways:
\begin{itemize}
  \item agreement across samples \citep{wang2023selfconsistency};
  \item faithfulness and coherence of chains of thought
    \citep{turpin2023say,lanham2023measuring};
  \item self-knowledge \citep{kadavath2022know}.
\end{itemize}
We use the six-dimension decomposition of RQEval \citep{rqeval2026}:
\begin{itemize}
  \item answer accuracy (CQ);
  \item consistency (CS);
  \item robustness (RS);
  \item logical coherence (LS), computed with an NLI model \citep{he2021deberta};
  \item efficiency (ES);
  \item stability (SS), computed with BERTScore \citep{zhang2020bertscore}.
\end{itemize}
We test whether these dimensions warn; we do not modify them.

\paragraph{Drift detection.}
Classic stream detectors include:
\begin{itemize}
  \item CUSUM \citep{page1954continuous};
  \item Page--Hinkley \citep{hinkley1971inference};
  \item EWMA charts \citep{roberts1959control};
  \item DDM \citep{gama2004learning};
  \item ADWIN \citep{bifet2007learning}.
\end{itemize}
Surveys of concept drift include \citet{gama2014survey} and
\citet{lu2019learning}. Shift detection for deployed models is studied by
\citet{rabanser2019failing}, and behavioural drift of commercial LLMs over time
by \citet{chen2024chatgpt}. We apply these detectors to per-turn
reasoning-quality streams under calibrated false-alarm control.

\paragraph{Prospective analysis plans.}
Model evaluations are prone to analytic flexibility in prompts, parsers and
thresholds. Fixing hypotheses and analyses in advance \citep{nosek2018preregistration}
limits it, so we froze each study's protocol before its data, and we label every
post-hoc analysis as such.
